\SourceLeaf{202}
\textit{Ex.} - Desiderantur omnes transformationes formae $(129,92,65)$ in formam $(42,59,81)$. Has improprie tantum aequivalentes esse in art.~195 invenimus et in art. seq. transformationem impropriam illius in hanc eruimus $-47,-56,73,87$. Quamobrem omnes transformationes formae $(129,92,65)$ in $(42,59,81)$ exhibebuntur per formulam
\[
-(47t+421u),\quad -(56t+503u).\quad 73t+653u,\quad 87t+780u
\]
ubi $t$, $u$ sunt indefinite omnes numeri aequationi $tt-79uu=1$ satisfacientes; hi vero exhibentur per formulas
\[
\begin{aligned}
\pm t&=\tfrac12\bigl((80+9\sqrt{79})^e+(80-9\sqrt{79})^e\bigr)\\
\pm u&=\tfrac{1}{2\sqrt{79}}\bigl((80+9\sqrt{79})^e-(80-9\sqrt{79})^e\bigr)
\end{aligned}
\]
ubi pro $e$ omnes numeri integri non negativi sunt accipiendi.

\Needspace{8\baselineskip}
\Article{204}
Perspicuum est, formulam generalem omnes transformationes exhibentem eo \textit{simpliciorem} evadere, quo simplicior fuerit transformatio initialis ex qua formula est deducta. Iam quum arbitrarium sit, a qua transformatione proficiscamur, saepenumero formula generalis simplicior reddi potest, si ex formula primo inventa transformatio simplicior deducitur tribuendo ipsis $t$, $u$ valores determinatos, et tunc ex hac alia formula componitur. Ita \textit{e.~g.} positis in formula in ex. art. praec. inventa, $t=80$, $u=-9$, prodit transformatio simplicior quam ea a qua profecti eramus, scilicet $29,47,-37,-60$, unde deducitur formula generalis $29t-263u$, $47t-424u$, $-37t+337u$, $-60t+543u$. Quando itaque per praecepta praecedentia formula generalis eruta est, tentari poterit, annon, tribuendo ipsis $t$, $u$ valores determinatos $\pm t',\pm u';\ \pm t'',\pm u''$ etc. transformatio obtineatur simplicior quam ea ex qua formula deducta fuit, in quo casu ex illa transformatione formula simplicior derivari poterit. --- Ceterum in diiudicanda simplicitate aliquid arbitrarii remanet, quod si operae pretium esset ad normam fixam revocare, nec non in progressione $t',u';\ t'',u''$ etc. \textit{limites} assignare possemus, ultra quos transformationes continuo minus simplices prodeant, ita ut ultra progredi opus non sit sed intra illos tentamen instituisse sufficiat: attamen quum plerumque per methodos a nobis praescriptas transformatio simplicissima vel sta%
\SourceLeaf{203}%
tim vel adhibitis pro $t$, $u$ valoribus $\pm t',\pm u'$ prodire soleat, hanc disquisitionem brevitatis gratia supprimimus.

\Needspace{8\baselineskip}
\Article{205}
\textsc{Problema.} \textit{Invenire omnes repraesentationes numeri dati $M$ per formulam datam $axx+2bxy+cyy$, cuius determinans positivus non-quadratus $=D$.}

\textit{Sol.} Primo observamus, investigationem repraesentationum per valores ipsorum $x$, $y$ inter se non primos, hic prorsus eodem modo, ut supra (art.~181) pro formis determinantis negativi, ad eum casum reduci posse, ubi repraesentationes per valores indeterminatarum inter se primos quaeruntur, quod igitur hic repetere superfluum foret. Ad possibilitatem repraesentationum per valores ipsorum $x,y$ inter se primos autem requiritur, ut $D$ sit residuum quadraticum ipsius $M$, et si omnes valores expressionis $\sqrt D\;(\mathrm{mod.}\,M)$ sunt $N,-N,N',-N',N'',-N''$ etc. (quos ita accipere licet ut nullus sit $>\tfrac12M$), quaevis repraesentatio numeri $M$ per formam propositam ad aliquem horum valorum pertinebit. Ante omnia itaque valores illi erui debebunt; tunc repraesentationes ad singulos pertinentes deinceps investigari. Repraesentationes ad valorem $N$ pertinentes non dabuntur, nisi formae $(a,b,c)$ et $(M,N,\frac{NN-D}{M})$ proprie aequivalentes sunt; si vero sunt, quaeratur transformatio aliqua propria prioris in posteriorem, quae sit $\alpha,\beta,\gamma,\delta$. Tum habebitur repraesentatio numeri $M$ per formam $(a,b,c)$ ad valorem $N$ pertinens haec: $x=\alpha$, $y=\gamma$, omnesque repraesentationes ad hunc valorem pertinentes exhibebuntur per formulam
\[
x=\tfrac1m\bigl(\alpha t-(\alpha b+\gamma c)u\bigr),\qquad
y=\tfrac1m\bigl(\gamma t+(\alpha a+\gamma b)u\bigr)
\]
designante $m$ divisorem communem maximum numerorum $a,2b,c$; et $t$, $u$ indefinite omnes numeros aequationi $tt-Duu=mm$ satisfacientes. --- Ceterum manifestum est, formulam hanc generalem eo simpliciorem evadere, quo simplicior sit transformatio $\alpha,\beta,\gamma,\delta$ ex qua deducta est; quare haud inutile erit, transformationem simplicissimam formae $(a,b,c)$ in $(M,N,\frac{NN-D}{M})$ secundum art. praec. antea eruere, et ex hac formulam deducere. --- Prorsus eodem modo repraesentationes ad valores reliquos $-N,N',-N'$ etc. pertinentes (si quae dantur) per formulas generales exhiberi possunt.

\textit{Ex.} Quaerantur omnes repraesentationes numeri $585$ per formulam
\SourceLeaf{204}
$42xx+62xy+21yy$. Quod ad repraesentationes per valores ipsorum $x$, $y$ inter se non primos pertinet, statim patet alias huius generis dari non posse, quam in quibus divisor communis maximus ipsorum $x$, $y$ sit $3$: quum $585$ per unicum quadratum $9$ divisibilis sit. Quando itaque omnes repraesentationes numeri $\frac{585}{9}$ \textit{i.~e.} $65$ per formam $42x'x'+62x'y'+21y'y'$ inventae sunt, in quibus $x'$ ad $y'$ primus; omnes repraesentationes numeri $585$ per formam $42xx+62xy+21yy$, in quibus $x$ ad $y$ non primus, ex illis derivabuntur ponendo $x=3x'$, $y=3y'$. Valores expressionis $\sqrt{79}\;(\mathrm{mod.}\,65)$ sunt $\pm12,\pm27$. Repraesentatio numeri $65$ ad valorem $-12$ pertinens invenitur $x'=2$, $y'=-1$; quocirca omnes repraesentationes ipsius $65$ ad hunc valorem pertinentes exhibebuntur per formulam $x'=2t-41u$, $y'=-t+53u$, adeoque omnes repraesentationes ipsius $585$ \textit{hinc} oriundae per formulam $x=6t-123u$, $y=-3t+159u$. Simili modo invenitur formula generalis omnes repraesentationes numeri $65$ ad valorem $+12$ pertinentes exhibens $x'=22t-199u$, $y'=-23t+211u$; et formula omnes repraesentationes numeri $585$ hinc oriundas complectens $x=66t-597u$, $y'=-69t+633u$. Ad valores $+27$ et $-27$ autem nulla repraesentatio numeri $65$ pertinet. --- Ut repraesentationes numeri $585$ per valores ipsorum $x$, $y$ inter se primos inveniantur, primo valores expressionis $\sqrt{79}\;(\mathrm{mod.}\,585)$ eruere oportet, qui sunt $\pm77,\pm103,\pm157,\pm248$. Ad valores $\pm77,\pm103,\pm248$ invenitur nullam repraesentationem pertinere; ad valorem $-157$ autem pertinet repraesentatio $x=3$, $y=1$, unde deducitur formula generalis omnes repraesentationes ad hunc valorem pertinentes exhibens $x=3t-114u$, $y=t+157u$; similiterque invenitur repraesentatio ad $+157$ pertinens $x=83$, $y=-87$, et formula in qua omnes similes sunt contentae $x=83t-746u$, $y=-87t+789u$. Habentur itaque quatuor formulae generales, sub quibus omnes repraesentationes numeri $585$ per formam $42xx+62xy+21yy$ contentae sunt
\[
\begin{aligned}
x&=6t-123u&\qquad y&=-3t+159u\\
x&=66t-597u& y&=-69t+633u\\
x&=3t-114u& y&=t+157u\\
x&=83t-746u& y&=-87t+789u
\end{aligned}
\]
ubi $t,u$ indefinite omnes numeros integros denotant, qui aequationi $tt-79uu=1$ satisfaciunt.

Applicationibus specialibus disquisitionum praecedentium de formis deter%
\SourceLeaf{205}%
minantis positivi non-quadrati brevitatis gratia non immoramur, quippe quas simili modo ut artt.~176, 182 quisque, sine negotio, proprio marte instituere poterit, statimque ad formas determinantis positivi quadrati, quae solae adhuc supersunt, properamus.

\Needspace{12\baselineskip}
\begin{center}
\textit{De formis determinantis quadrati.}
\end{center}

\Needspace{8\baselineskip}
\Article{206}
\textsc{Problema.} \textit{Proposita forma $(a,b,c)$ determinantis quadrati $hh$, designante $h$ ipsius radicem positivam, invenire formam $(A,B,C)$ illi proprie aequivalentem, in qua $A$ iaceat inter limites $0$ et $2h-1$ incl., $B$ sit $=h$, $C=0$.}

\textit{Sol.} I. Quoniam $hh=bb-ac$, erit $(h-b):a=c:-(h+b)$. Sit huic rationi aequalis ratio $\beta:\delta$, ita ut $\beta$ ad $\delta$ sit primus, determinenturque $\alpha$, $\gamma$ ita ut sit $\alpha\delta-\beta\gamma=1$, quae fieri poterunt. Per substitutionem $\alpha,\beta,\gamma,\delta$ transeat forma $(a,b,c)$ in $(a',b',c')$, quae igitur illi proprie aequivalens erit. Habebitur autem
\[
\begin{aligned}
b'&=a\alpha\beta+b(\alpha\delta+\beta\gamma)+c\gamma\delta\\
&=(h-b)\alpha\delta+b(\alpha\delta+\beta\gamma)-(h+b)\beta\gamma\\
&=h(\alpha\delta-\beta\gamma)=h\\
c'&=a\beta\beta+2b\beta\delta+c\delta\delta\\
&=(h-b)\beta\delta+2b\beta\delta-(h+b)\beta\delta=0
\end{aligned}
\]
Quodsi itaque insuper $a'$ inter limites $0$ et $2h-1$ iam est situs, forma $(a',b',c')$ omnibus conditionibus satisfaciet.

II. Si vero $a'$ extra limites $0$ et $2h-1$ iacet, sit $A$ residuum minimum positivum ipsius $a'$ secundum modulum $2h$, quod manifesto inter hos limites situm erit, ponaturque $A-a'=2hk$. Tum forma $(a',b',c')$ \textit{i.~e.} $(a',h,0)$ per substitutionem $1,0,k,1$ transibit in formam $(A,h,0)$, quae formis $(a',b',c')$, $(a,b,c)$ proprie aequivalens erit omnibusque conditionibus satisfaciet. --- Ceterum perspicuum est, formam $(a,b,c)$ transire in formam $(A,h,0)$ per substitutionem $\alpha+\beta k,\beta,\gamma+\delta k,\delta$.

\textit{Ex.} Proposita sit forma $(27,15,8)$ cuius determinans $=9$. Hic $h=3$; rationibus $-12:27=8:-18$ in numeris minimis aequalis est ratio $4:-9$. Positis itaque $\beta=4$, $\delta=-9$, $\alpha=-1$, $\gamma=2$, forma $(a',b',c')$ fit $(-1,3,0)$. quae transit in formam $(5,3,0)$ per substitutionem $1,0,1,1$. Haec igitur est
\SourceLeaf{206}
forma quaesita, transitque in eam proposita per substitutionem propriam $3,4,-7,-9$.

Tales formas $(A,B,C)$, in quibus $C=0$, $B=h$, $A$ inter limites $0$ et $2h-1$ situs, \textit{formas reductas} vocabimus, quae igitur a formis reductis determinantis negativi, vel positivi non-quadrati, probe sunt distinguendae.

\Needspace{8\baselineskip}
\Article{207}
\textsc{Theorema.} \textit{Duae formae reductae $(a,h,0)$, $(a',h,0)$, non identicae proprie aequivalentes esse non possunt.}

\textit{Dem.} Si enim proprie aequivalere supponuntur, transeat prior in posteriorem per substitutionem propriam $\alpha,\beta,\gamma,\delta$, habebunturque quatuor aequationes:
\begin{align*}
a\alpha\alpha+2h\alpha\gamma&=a'\tag*{. . . . . . . . . [1]}\\
a\alpha\beta+h(\alpha\delta+\beta\gamma)&=h\tag*{. . . . . . . [2]}\\
a\beta\beta+2h\beta\delta&=0\tag*{. . . . . . . . . . [3]}\\
\alpha\delta-\beta\gamma&=1\tag*{. . . . . . . . . . . [4]}
\end{align*}
Multiplicando aequationem secundam per $\beta$, tertiam per $\alpha$ et subtrahendo fit $-h(\alpha\delta-\beta\gamma)\beta=\beta h$, sive, propter [4], $-\beta h=\beta h$, unde necessario $\beta=0$. Quare ex [4], $\alpha\delta=1$, et $\alpha=\pm1$. Hinc ex [1], $a\pm2\gamma h=a'$. quae aequatio consistere nequit, nisi $\gamma=0$ (quoniam tum $a$ tum $a'$ per hyp. inter $0$ et $2h-1$ iacent) \textit{i.~e.} nisi $a=a'$, sive formae $(a,h,0)$, $(a',h,0)$ identicae, contra hyp.

Hinc sequentia problemata, quae pro determinantibus non-quadratis multo maiorem difficultatem facessebant, nullo negotio solvi poterunt.

I. \textit{Propositis duabus formis $F$, $F'$ eiusdem determinantis quadrati, investigare an proprie aequivaleant.} Quaerantur duae formae reductae formis $F$, $F'$ resp. proprie aequivalentes; quae si identicae sunt, propositae proprie aequivalentes erunt, sin minus, non erunt.

II. \textit{Iisdem positis investigare an improprie aequivaleant.} Sit forma alterutri propositarum \textit{e.~g.} formae $F$ opposita, $G$; quae si formae $F'$ proprie aequivalet, $F$ et $F'$ improprie aequivalebunt, et contra.

\Needspace{8\baselineskip}
\Article{208}
\textsc{Problema.} \textit{Propositis duabus formis $F$, $F'$ determinantis $hh$ proprie aequivalentibus: invenire transformationem propriam alterius in alteram.}

\SourceLeaf{207}
\textit{Sol.} Formae $F$ proprie aequivaleat forma reducta $\Phi$, quae itaque per hyp. etiam formae $F'$ proprie aequivalebit. Quaeratur per art.~206 transformatio propria formae $F$ in $\Phi$, quae sit $\alpha,\beta,\gamma,\delta$; nec non transformatio propria formae $F'$ in $\Phi$, quae sit $\alpha',\beta',\gamma',\delta'$. Tunc $\Phi$ transformabitur in $F'$ per substitutionem propriam $\delta',-\beta',-\gamma',\alpha'$ et hinc $F$ in $F'$ per substitutionem propriam
\[
\alpha\delta'-\beta\gamma',\quad\beta\alpha'-\alpha\beta',\quad
\gamma\delta'-\delta\gamma',\quad\delta\alpha'-\gamma\beta'
\]

Operae pretium est, aliam formulam pro hac transformatione formae $F$ in $F'$ evolvere, ad quam formam reductam $\Phi$ ipsam novisse ne opus quidem sit. Ponamus formam
\[
F\text{ esse }(a,b,c),\quad F'=(a',b',c'),\quad\Phi=(A,h,0)
\]
Quoniam rationibus $h-b:a$ vel $c:-(h+b)$ in numeris minimis aequalis est ratio $\beta:\delta$, facile perspicitur $\frac{h-b}{\beta}=\frac a\delta$ fore \textit{integrum}, qui sit $f$; nec non $\frac c\beta=\frac{-h-b}{\delta}$ integrum fore qui ponatur $=g$. Habebitur autem
\[
\begin{gathered}A=a\alpha\alpha+2b\alpha\gamma+c\gamma\gamma\quad\text{adeoque}\\
\beta A=a\alpha\alpha\beta+2b\alpha\beta\gamma+c\beta\gamma\gamma\end{gathered}
\]
sive (substitutis pro $a\beta$, $\delta(h-b)$, pro $c$, $\beta g$)
\[
\beta A=\alpha\alpha\delta h+b(2\beta\gamma-\alpha\delta)\alpha+\beta\beta\gamma\gamma g
\]
sive (propter $b=-h-\delta g$)
\[
\beta A=2\alpha(\alpha\delta-\beta\gamma)h+(\alpha\delta-\beta\gamma)^2g=2\alpha h+g
\]
Simili modo
\[
\begin{aligned}
\delta A&=a\alpha\alpha\delta+2b\alpha\gamma\delta+c\gamma\gamma\delta\\
&=\alpha\alpha\delta\delta f+b(2\alpha\delta-\beta\gamma)\gamma-\beta\gamma\gamma h\\
&=(\alpha\delta-\beta\gamma)^2f+2\gamma(\alpha\delta-\beta\gamma)h=2\gamma h+f
\end{aligned}
\]
Quare
\[
\alpha=\tfrac{\beta A-g}{2h},\qquad\gamma=\tfrac{\delta A-f}{2h}
\]
Prorsus eodem modo positis
\[
\frac{h-b'}{\beta'}=\frac{a'}{\delta'}=f',\qquad
\frac{c'}{\beta'}=\frac{-h-b'}{\delta'}=g'
\]
fit
\[
\alpha'=\tfrac{\beta' A-g'}{2h},\qquad\gamma'=\tfrac{\delta' A-f'}{2h}
\]
\SourceLeaf{208}
Quibus valoribus ipsorum $\alpha,\gamma,\alpha',\gamma'$ in formula modo tradita pro transformatione formae $F$ in $F'$ substitutis, transit in hanc:
\[
\tfrac{\beta f'-\delta'g}{2h},\quad
\tfrac{\beta'g-\beta g'}{2h},\quad
\tfrac{\delta f'-\delta'f}{2h},\quad
\tfrac{\beta'f-\delta g'}{2h}
\]
ex qua $A$ omnino abiit.

Si duae formae $F$, $F'$ improprie aequivalentes proponuntur, et transformatio impropria alterius in alteram quaeritur, sit forma $G$ opposita formae $F$, et transformatio propria formae $G$ in $F'$ haec $\alpha,\beta,\gamma.\ \delta$. Tunc manifestum est, $\alpha,\beta,-\gamma,-\delta$ fore transformationem impropriam formae $F$ in $F'$.

Denique patet, si formae propositae et proprie et improprie aequivalentes sint, hoc modo inveniri posse transformationes duas alteram propriam alteram impropriam.

\Needspace{8\baselineskip}
\Article{209}
Nihil itaque iam superest quam ut ex una transformatione omnes reliquas similes deducere doceamus. Hoc vero pendet a solutione aequationis indeterminatae $tt-hhuu=mm$, designante $m$ divisorem communem maximum numerorum $a,2b,c$, si $(a,b,c)$ est alterutra formarum aequivalentium. Sed haec aequatio semper duobus tantum modis solvi potest, nempe ponendo aut $t=m$, $u=0$, aut $t=-m$, $u=0$. Ponamus enim dari adhuc aliam solutionem $t=T$, $u=U$, ita ut $U$ non $=0$. Quia $mm$ ipsum $4hh$ certo metitur, erit $\frac{4TT}{mm}=\frac{4hhUU}{mm}+4$, atque tum $\frac{4TT}{mm}$ tum $\frac{4hhUU}{mm}$ quadrata integra. Sed nullo negotio perspicitur, numerum $4$ duorum quadratorum integrorum differentiam esse non posse, nisi quadratum minus sit $0$ \textit{i.~e.} $U=0$, contra hyp. --- Si itaque forma $F$ in formam $F'$ per substitutionem $\alpha,\beta,\gamma,\delta$ transit, alia transformatio huic similis non dabitur praeter transformationem $-\alpha,-\beta,-\gamma,-\delta$. Quare si duae formae aut proprie tantum, aut improprie tantum aequivalent, \textit{duae} tantum transformationes dabuntur; si vero tum proprie tum improprie, \textit{quatuor}, nempe duae propriae duaeque impropriae.

\Needspace{8\baselineskip}
\Article{210}
\textsc{Theorema.} \textit{Si duae formae reductae $(a,h,0)$, $(a',h,0)$ improprie sunt aequivalentes, erit $aa'\equiv mm\;(\mathit{mod.}\,2mh)$, designante $m$ divisorem communem maximum numerorum $a,2h$, vel $a',2h$; et vice versa, si $a,2h;\ a',2h$ eundem divisorem}
\SourceLeaf{209}
\textit{communem maximum $m$ habent, atque est $aa'\equiv mm\;(\mathit{mod.}\,2mh)$, formae $(a,h,0)$, $(a',h,0)$ improprie aequivalentes erunt.}

\textit{Dem.} I. Transeat forma $(a,h,0)$ in formam $(a',h,0)$ per substitutionem impropriam $\alpha,\beta,\gamma,\delta$ ita ut habeantur quatuor aequationes
\begin{align*}
a\alpha\alpha+2h\alpha\gamma&=a'\tag*{. . . . . . . . . . . . [1]}\\
a\alpha\beta+h(\alpha\delta+\beta\gamma)&=h\tag*{. . . . . . . . . . [2]}\\
a\beta\beta+2h\beta\delta&=0\tag*{. . . . . . . . . . . . [3]}\\
\alpha\delta-\beta\gamma&=-1\tag*{. . . . . . . . . . . . . [4]}
\end{align*}
Hinc sequitur, multiplicando [4] per $h$ et subtrahendo a [2], quod ita exprimimus $[2]-h[4]$,
\begin{equation*}
(a\alpha+2h\gamma)\beta=2h\tag*{. . . . . . . . . . . . [5]}
\end{equation*}
similiter ex $\gamma\delta[2]-\gamma\gamma[3]-(a+a\beta\gamma+h\gamma\delta)[4]$ deletis quae sese destruunt
\begin{equation*}
-a\alpha\delta=a+2h\gamma\delta,\quad\text{sive}\quad -(a\alpha+2h\gamma)\delta=a\tag*{. . [6]}
\end{equation*}
denique ex $a[1]\ldots a\alpha(a\alpha+2h\gamma)=aa'$, sive
\[
(a\alpha+2h\gamma)^2-aa'=2h\gamma(a\alpha+2h\gamma)
\]
sive
\begin{equation*}
(a\alpha+2h\gamma)^2\equiv aa'\;(\mathrm{mod.}\,2h(a\alpha+2h\gamma))\tag*{. . . . . [7]}
\end{equation*}
Iam ex [5] et [6] sequitur $a\alpha+2h\gamma$ metiri ipsos $2h$ et $a$, adeoque etiam ipsum $m$, qui est divisor communis maximus ipsorum $a,2h$; manifesto autem $m$ metietur etiam ipsum $a\alpha+2h\gamma$; quare necessario $a\alpha+2h\gamma$ erit aut $=+m$ aut $=-m$. Hinc statim sequitur ex [7], $mm\equiv aa'\;(\mathrm{mod.}\,2mh)$ \textit{Q.~E.~P.}

II. Si $a,2h;\ a',2h$ eundem divisorem communem maximum $m$ habent, insuperque est $aa'\equiv mm\;(\mathrm{mod.}\,2mh)$, $\frac a m$, $\frac{2h}{m}$, $\frac{a'}{m}$, $\frac{aa'-mm}{2mh}$ erunt integri. Facile vero confirmatur, formam $(a,h,0)$ transire in $(a',h,0)$ per substitutionem $-\frac{a'}{m},-\frac{2h}{m},\frac{aa'-mm}{2mh},\frac a m$; nec non hanc transformationem esse impropriam. Quare formae illae erunt improprie aequivalentes. \textit{Q.~E.~S.}

Hinc etiam statim diiudicari potest, an forma aliqua reducta data $(a,h,0)$ sibi ipsi improprie aequivalens sit. Scilicet designato divisore communi maximo numerorum $a,2h$ per $m$, esse debebit $aa\equiv mm\;(\mathrm{mod.}\,2mh)$.

\SourceLeaf{210}
\Needspace{8\baselineskip}
\Article{211}
Omnes formae reductae determinantis dati $hh$ obtinentur, si in forma indefinita $(A,h,0)$ pro $A$ omnes numeri a $0$ usque ad $2h-1$ incl substituuntur, quarum itaque multitudo erit $2h$. Perspicuum est, omnes formas determinantis $hh$ in totidem \textit{classes} distribui posse, hasque iisdem proprietatibus praeditas fore quas supra (artt.~175, 195) pro classibus formarum determinantis negativi, et positivi non-quadrati attigimus. Ita omnes formae determinantis $25$ in decem classes distribuentur, quae per formas reductas in singulis contentas distingui poterunt. Hae formae reductae sunt: $(0,5,0)$, $(1,5,0)$, $(2,5,0)$, $(5,5,0)$, $(8,5,0)$, $(9,5,0)$, quae sibi ipsae simul improprie aequivalent; $(3,5,0)$ cui improprie aequivalet $(7,5,0)$; $(4,5,0)$ cui improprie aequivalet $(6,5,0)$.

\Needspace{8\baselineskip}
\Article{212}
\textsc{Problema.} \textit{Invenire omnes repraesentationes numeri dati $M$ per formam datam $axx+2bxy+cyy$ determinantis $hh$.}

Solutio huius problematis ex principiis art.~168 prorsus eodem modo peti potest, ut supra (artt.~180, 181, 205) pro formis determinantis negativi et positivi non-quadrati ostendimus; quod, quum nulli difficultati sit obnoxium, hic repetere superfluum esset. Contra haud abs re erit, solutionem ex alio principio quod casui praesenti proprium est deducere.

Positis ut artt.~206, 208
\[
\begin{gathered}
h-b:a=c:-(h+b)=\beta:\delta\\
\frac{h-b}{\beta}=\frac a\delta=f;\quad
\frac c\beta=\frac{-h-b}{\delta}=g
\end{gathered}
\]
nullo negotio probatur, formam propositam esse productum ex factoribus $\delta x-\beta y$ et $fx-gy$. Unde manifestum est, quamvis repraesentationem numeri $M$ per formam propositam praebere resolutionem numeri $M$ in binos factores. Si itaque omnes divisores numeri $M$ sunt $d,d',d''$ etc. (inclusis etiam $1$, et $M$, et singulis \textit{bis} sumtis puta tum positive tum negative), patet omnes repraesentationes numeri $M$ obtineri, si successive ponatur
\[
\begin{aligned}
\delta x-\beta y&=d,&\quad fx-gy&=\frac M d\\
\delta x-\beta y&=d',& fx-gy&=\frac M{d'}\quad\text{etc.}
\end{aligned}
\]
valores ipsorum $x$, $y$ hinc evolvantur, eaeque repraesentationes eiiciantur ubi $x$
\SourceLeaf{211}
aut $y$ valores fractos obtinent. Manifesto vero ex duabus primis aequationibus sequitur
\[
x=\tfrac{\beta M-gdd}{(\beta f-\delta g)d}\quad\text{et}\quad
y=\tfrac{\delta M-fdd}{(\beta f-\delta g)d}
\]
quos valores semper \textit{determinatos} fore inde manifestum quod $\beta f-\delta g=2h$, adeoque numerator certo non $=0$. --- Ceterum ex eodem principio, puta resolubilitate cuiusvis formae determinantis quadrati in binos factores, etiam reliqua problemata solvi potuissent: sed methodo ei quam supra pro formis determinantis non-quadrati tradidimus analoga etiam hic uti maluimus.

\textit{Ex.} Quaeruntur omnes repraesentationes numeri $12$ per formam $3xx+4xy-7yy$. Haec resolvitur in factores $x-y$ et $3x+7y$. Omnes divisores numeri $12$ sunt $\pm1,2,3,4,6,12$. Positis $x-y=1$, $3x+7y=12$ fit $x=\frac{19}{10}$, $y=\frac9{10}$, qui valores tamquam fracti sunt reiiciendi. Eodem modo ex divisoribus $-1,\pm3,\pm4,\pm6,\pm12$ valores inutiles obtinentur; ex divisore $+2$ vero obtinentur valores $x=2,y=0$, et ex divisore $-2$ hi $x=-2$, $y=0$; praeter has duas repraesentationes igitur aliae non dantur.

Methodus haec adhiberi nequit, si $M=0$. In hoc casu manifestum est omnes valores ipsorum $x,y$ aut aequationi $\delta x-\beta y=0$, aut huic $fx-gy=0$ satisfacere debere. Omnes autem solutiones aequationis prioris continentur in formula $x=\beta z$, $y=\delta z$, designante $z$ indefinite numerum integrum quemcunque (siquidem uti supponitur $\beta$, $\delta$ inter se primi sunt); similiterque ponendo divisorem communem maximum numerorum $f,g$, $=m$, omnes solutiones aequationis posterioris exhibebuntur per formulam $x=\frac{gz}{m},y=\frac{hz}{m}$. Quare hae duae formulae generales omnes repraesentationes numeri $M$ in hoc casu complectentur.

\begin{center}
\(\begin{matrix} *&\qquad&*\\[-0.8ex]&*&\end{matrix}\)
\end{center}

In praecedentibus omnia quae ad cognoscendam aequivalentiam et ad inveniendas omnes transformationes formarum nec non ad repraesentationes omnes numerorum datorum per formas datas indagandas pertinent, ita sunt explicata, ut nihil amplius desiderari posse videatur. Superest itaque tantummodo, ut propositis duabus formis quae propter \textit{determinantium inaequalitatem} aequivalentes esse nequeunt, diiudicare doceamus, annon altera sub altera contenta sit. et in hoc casu omnes transformationes illius in hanc invenire.
